%% LyX 2.4.4 created this file.  For more info, see https://www.lyx.org/.
%% Do not edit unless you really know what you are doing.
\documentclass[english]{article}
\usepackage[T1]{fontenc}
\usepackage[latin9]{inputenc}
\usepackage{enumitem}
\usepackage{amsmath}
\usepackage{amsthm}
\usepackage{geometry}
\geometry{verbose,tmargin=1in,bmargin=1in,lmargin=1in,rmargin=1in}

\makeatletter
%%%%%%%%%%%%%%%%%%%%%%%%%%%%%% Textclass specific LaTeX commands.
\numberwithin{equation}{section}
\numberwithin{figure}{section}
\newlength{\lyxlabelwidth}      % auxiliary length 

%%%%%%%%%%%%%%%%%%%%%%%%%%%%%% User specified LaTeX commands.
\usepackage{fancyhdr}
\usepackage{lastpage}
\pagestyle{fancy}
\usepackage{enumitem}
\usepackage{ifthen}
\everymath=\expandafter{\the\everymath\displaystyle}
%\DeclareMathSizes{10}{10}{10}{10}
\usepackage{multicol}

\usepackage{tikz}
\usepackage{tikz-3dplot}
\usetikzlibrary{calc,arrows}

\usetikzlibrary{patterns}
\usetikzlibrary{plotmarks}
\usepackage{pgfplots}
\pgfplotsset{compat=newest}

\usepackage{calc}
\usepackage[nomessages]{fp}

\usepackage{datetime}
% Added by lyx2lyx
\setlength{\parskip}{\smallskipamount}
\setlength{\parindent}{0pt}

\makeatother

\usepackage{babel}
\begin{document}
\renewcommand{\author}{Dr.\,James Ashe}

\newcommand{\classNum}{336}

\newcommand{\classSec}{A}

\newcommand{\examType}{HW 2.1 Solutions}

\renewcommand{\date}{\formatdate{30}{9}{2013}}

%%First page's header%%%%%%%%%%%%%%%%%%%%%%
\fancypagestyle{firststyle} %%header settings for first pagr
{
	\fancyhead[L]{MTH \classNum{}(\classSec{}) \author{}\\
	\textbf{\examType{}} ~\date{}}
	\fancyhead[C]{}
	\fancyhead[R]{\textbf{Name:\hspace{4cm}}}
	\setlength{\headsep}{14pt}
}
%%%%%%%%%%%%%%%%%%%%%%%%%%%%%%%%%%%%%%%%%%

%%Next page's header%%%%%%%%%%%%%%%%%%%%%%%
\setlist{nolistsep}
\lhead{\classNum{}(\classSec{})} 
\chead{\textbf{Name:}} 
\cfoot{\thepage{}~of~\pageref{LastPage}} 
\setlength{\headsep}{5pt} 
%%%%%%%%%%%%%%%%%%%%%%%%%%%%%%%%%%%%%%%%%%%

%%Various settings; see the preamble for other settings%%%%%
%\setenumerate[0]{leftmargin=12pt,itemindent=0pt,labelwidth=0pt}
\setlist{topsep=.5pt}
\renewcommand{\labelenumi}{\textbf{\arabic{enumi}.}}
\renewcommand{\labelenumii}{\textbf{\alph{enumii}.}}
\thispagestyle{firststyle}
%%%%%%%%%%%%%%%%%%%%%%%%%%%%%%%%%%%%%%%%%%

\newcommand{\xmax}{0}
\newcommand{\xmin}{-\xmax}
\newcommand{\xstep}{0}
\newcommand{\ymax}{0}
\newcommand{\ymin}{-\ymax}
\newcommand{\ystep}{0} 
\newcommand{\dspace}{\vspace{1cm}}
\newcommand{\xa}{0}
\newcommand{\xb}{0}
\newcommand{\ya}{1}
\newcommand{\yb}{1}
\begin{enumerate}[resume]
\item \vspace{0cm}

\begin{enumerate}[resume]
\item $A+B=\begin{bmatrix}\begin{array}{rr}
1+2 & 2+1\\
3+4 & 4+3
\end{array}\end{bmatrix}=\begin{bmatrix}\begin{array}{rr}
3 & 3\\
7 & 7
\end{array}\end{bmatrix}$
\item $2A=\begin{bmatrix}\begin{array}{rr}
2\cdot1 & 2\cdot2\\
2\cdot3 & 2\cdot4
\end{array}\end{bmatrix}=\begin{bmatrix}\begin{array}{rr}
2 & 4\\
6 & 8
\end{array}\end{bmatrix}$; so $2A-B=\begin{bmatrix}\begin{array}{rr}
2-2 & 4-1\\
6-4 & 8-3
\end{array}\end{bmatrix}=\begin{bmatrix}\begin{array}{rr}
0 & 3\\
2 & 5
\end{array}\end{bmatrix}$
\item [\textbf{d.}]\setcounter{enumii}{5}$C+D$ is undefined. Since $C$
and $D$ have different dimensions.
\item $BA=\begin{bmatrix}\begin{array}{rr}
2\cdot1+1\cdot3 & 2\cdot2+1\cdot4\\
4\cdot1+3\cdot3 & 4\cdot2+3\cdot4
\end{array}\end{bmatrix}=\begin{bmatrix}\begin{array}{rr}
5 & 8\\
13 & 20
\end{array}\end{bmatrix}$
\item $AC=\begin{bmatrix}\begin{array}{rrr}
1 & 4 & 5\\
3 & 10 & 11
\end{array}\end{bmatrix}$
\item $CA$ is undefined since $C$ has a different number of columns than
the rows of $A$.\vfill
\end{enumerate}
\item $AB=\mathbf{0}$ and $BA=\begin{bmatrix}5 & 5\\
-5 & 5
\end{bmatrix}$. This means that the vector rows of $A$ are orthogonal to the vector
columns of $B$, but the vector rows of $B$ are not orthogonal to
the vector columns of $A$. \vfill
\item Proposition 1.1 is on page 83. Let $A,B,C\in\mathcal{M}_{m\times n}$.
We'll denote the $ij^{\mbox{th}}$ entry in a matrix $A+B$ as $(A+B)_{ij}$.

\begin{enumerate}
\item Since real numbers are commutative, the $ij^{\mbox{th}}$ entry of
$A+B$ is $(A+B)_{ij}=a_{ij}+b_{ij}=b_{ij}+a_{ij}=(A+B)_{ij}$. So
the $A+B=B+A$.
\end{enumerate}
\begin{enumerate}[resume]
\item Since real numbers are associative, the $ij^{\mbox{th}}$ entry of
$\left((A+B)+C\right)_{ij}=(a_{ij}+b_{ij})+c_{ij}=a_{ij}+(b_{ij}+c_{ij})=\left(A+(B+C)\right)_{ij}$.
Thus $(A+B)+C=A+(B+C)$.
\item $a_{ij}+0=a_{ij}$. So then $\left(A+\mathbf{0}\right)_{ij}=a_{ij}+0=0+a_{ij}=a_{ij}=A_{ij}$.
Thus $A+\mathbf{0}=A$.\vfill
\end{enumerate}
\item [\textbf{6.}]Both of these statements are false. So we need only
provide counterexamples.

\begin{enumerate}
\item Let $A=\begin{bmatrix}1 & 0\\
1 & 0
\end{bmatrix}$ , $B=\begin{bmatrix}\begin{array}{rr}
1 & 0\\
0 & 0
\end{array}\end{bmatrix}$ and $C=\begin{bmatrix}\begin{array}{rr}
1 & 1\\
1 & 0
\end{array}\end{bmatrix}$. We then have $AB=\begin{bmatrix}\begin{array}{rr}
1 & 0\\
1 & 0
\end{array}\end{bmatrix}$ and $AC=\begin{bmatrix}\begin{array}{rr}
1 & 0\\
1 & 0
\end{array}\end{bmatrix}$. So $AB=AC$ but $A\neq C$ thus the claim is false. Note that this
shows that cancellation does not hold for matrix multiplication.  
\end{enumerate}
\begin{enumerate}[resume]
\item Let $A=\begin{bmatrix}\begin{array}{rr}
1 & 0\\
0 & 0
\end{array}\end{bmatrix}$. So then $A^{2}=\begin{bmatrix}\begin{array}{rr}
1 & 0\\
0 & 0
\end{array}\end{bmatrix}$ but $A\neq I_{2}$ and $A\neq\mathbf{0}$. \vfill
\end{enumerate}
\end{enumerate}

\end{document}
